Table~\ref{tab:therm} compares the geometries at the same gated operating point ($\Phi_M=5.0$~eV, $V_G=-2$~V). At $V_D=1$~V the edge-contacted device carries 3.7~A/cm and heats by 150~K; the 50-nm top contact carries 1.16~A/cm and heats by 24~K. Per unit power the edge device is also about twice as hot (41 vs 21~K\,cm/W), because its footprint on the oxide is a third of that of the top-contacted film. In the on state at $V_G=-3$~V (Fig.~\ref{fig:thermal}a) the edge device reaches about 490~K at 1~V, so its usable bias is thermally limited, not contact-limited. The two-dimensional maps in Fig.~\ref{fig:maps} show where the heat sits: in the top-contacted device the hottest point lies at the drain-side channel end on the contact side of the film, and the whole film stays within 24~K of ambient; in the edge-contacted device the film is 126--150~K hot everywhere and is coolest at the oxide interface, through which heat leaves to the substrate.

Table~\ref{tab:therm} uses the $S=1$ reference alignment, where $\Phi_M=5.0$~eV gives a hole barrier near zero. At the measured alignment (Pt, $V_G=-3$~V, $V_D=1$~V; Fig.~\ref{fig:measuredalign}a) self-heating depends on the pinning. Directly pinned contacts carry so little current that the edge device heats by only 0.2~K. With the oxide interlayer or the calculated ideal pinning ($\phi_{Bp}\approx0.52$~eV) the edge device heats by 24~K and the top-contacted one by 0.2~K. An unpinned Pt contact ($\phi_{Bp}=0.25$~eV) heats the edge device by 128~K (top: 4.4~K). Self-heating is therefore a concern only once the contacts are depinned. A three times larger $R_{\rm TBR}$ ($3\times10^{-8}$~m$^2$K/W) raises these rises by 58--73\%, and heat sinking through the metal lowers them by 22--23\%. Metal contacts are adiabatic in the baseline model, so the rises quoted are upper bounds. At the operating points of Table~\ref{tab:therm} the heat leaving the device equals $IV$ within 0.5\% (top) and 1.2--1.9\% (edge, 0.5 and 1~V). For the 150~K edge rise this is about 3~K, larger than the Peltier effect below, so temperatures are quoted to $\pm3$~K.

The individual heat terms matter at the 7--42\% level (Fig.~\ref{fig:thermal}b):
\begin{itemize}
\item Using bulk Joule heat only underestimates $\Delta T_{\max}$ by 7--8\% and misses 8--15\% of the dissipated power.
\item Measured out-of-plane conductivity, compared with isotropic $\kappa$, raises $\Delta T_{\max}$ by 8\% (edge) and 17\% (top).
\item Heat sinking through the metal contacts lowers $\Delta T_{\max}$ by 23\% (edge) and 42\% (top).
\item A ten times smaller thermal boundary resistance lowers it by 35--37\%.
\end{itemize}
Peltier heat changes $\Delta T_{\max}$ by only about 1~K but redistributes heat: it cools the hole-injecting contact and heats the collecting one (top, 1~V: source $+11.9$~K, drain $+21.6$~K with Peltier, against $+13.5$ and $+20.9$~K without). Temperature-dependent statistics change temperatures by less than 1\%. The lattice responds to a drain step within a few hundred picoseconds, with 63\% of the final rise at 0.20~ns (top) and 0.25~ns (edge) (Fig.~\ref{fig:thermal}c; this transient was simulated without tunnelling, so its amplitudes are smaller than in Table~\ref{tab:therm}). For pulses longer than a nanosecond the DC temperatures therefore apply.

\begin{table}[!htbp]
\centering
\caption{Gated operating points ($\Phi_M=5.0$~eV, $V_G=-2$~V, baseline thermal model). All quantities in a row refer to the same bias.}
\label{tab:therm}\small
\begin{tabular}{lcccccc}
\toprule
Geometry & $V_D$ (V) & $I_D$ (A/cm) & $IV$ (W/cm) & $\Delta T_{\max}$ (K) & $\Delta T_{\rm source}$ (K) & $\Delta T_{\rm drain}$ (K) \\
\midrule
Top, 50 nm & 0.5 & 0.618 & 0.309 & 5.9 & 3.5 & 5.6 \\
Top, 50 nm & 1.0 & 1.164 & 1.164 & 24.0 & 11.9 & 21.6 \\
Edge & 0.5 & 2.228 & 1.114 & 47.0 & 40.8 & 46.2 \\
Edge & 1.0 & 3.693 & 3.693 & 149.9 & 144.4 & 144.6 \\
\bottomrule
\end{tabular}
\end{table}
